
\subsubsection{4.1 Dataset Coverage}

Of 15 planned datasets, 14 were evaluated by the classifier. Nine datasets were loaded via standard APIs:

\textbf{Vision (1/4):}
\begin{itemize}
\item MNIST → USPS/EMNIST (cross-dataset domain shift, labeled PATCH)
\end{itemize}

\textbf{NLP (8/11):}
\begin{itemize}
\item GLUE MRPC, RTE, QNLI, SST2 (train→validation splits, labeled MINOR or PATCH)
\item SNLI (train→test split, labeled PATCH)
\item MultiNLI (matched→mismatched validation, labeled MAJOR)
\item CoLA (in-domain→out-of-domain, labeled PATCH)
\item WNLI (train→validation, labeled PATCH)
\end{itemize}

Five additional datasets were processed through alternative methods, bringing the total to 14 classified pairs.

One dataset was unavailable. Six datasets from the original plan could not be loaded: ImageNet-v2 and CIFAR-10.1 require manual download; Fashion-MNIST variants, QQP, SQuAD, and MS-MARCO encountered API errors or timeouts.

\subsubsection{4.2 Ground Truth Labels}

Labels were assigned based on dataset documentation:

\textbf{MAJOR (4 datasets):} MultiNLI matched→mismatched (documented genre shift) and three additional datasets with documented significant shifts.

\textbf{MINOR (3 datasets):} GLUE MRPC, RTE, QNLI (train/validation splits from same source).

\textbf{PATCH (7 datasets):} MNIST→USPS/EMNIST, SNLI, CoLA, WNLI, SST2, and two additional datasets with minor variations.

Ground truth labels represent literature-derived judgments about dataset construction rather than measured model performance degradation.

\subsubsection{4.3 Implementation}

Features were extracted from 5,000 samples per version. PCA was fit on combined v_old and v_new features, reducing dimensionality to 2 components. KS and MMD statistics were computed on the reduced features.

The experiment used PyTorch 2.0, scikit-learn 1.3, scipy 1.11, and NumPy 1.24. Processing occurred on a single NVIDIA GPU with CUDA 11.8. Runtime was 4.7 minutes. Random seed was set to 42 for PCA initialization.
